\chapter{Analysis Method}
I guess this is the part of the analysis that's not included in event selection.  
SO, how to get the cross section, like the main formula, 
and where all the numbers came from (especially the ones we haven't calculated yet).

%recap of invariant mass?

As explained in Section FIXME, 
the invariant mass of a set of decay-product particles is 
%CONCEPTUAL EXPLANATION.
essentially the mass of the original particle that decayed. 
%to form the end-products.  
It is defined mathematically as 

INV MASS FORMULA

where p momentum, E energy, etc. 
In our case, for a two-particle system, it can be written as 

MORE FORMULA (this one necessary??)

The invariant mass distribution for the electron pairs 
in the events surviving all selection cuts is shown in 
Fig.~\ref{fig:InvMass}.  
The peak around the Z mass (\~91 GeV) is clearly visible 
in both the Monte Carlo and data distributions.  

%yield plot first, since this is what the rest is based on.  

%%Figures: Electron-pair invariant mass for trigger objects/ID only/ID+isolation, Fig. \ref{fig:InvMassAfterEachStep}
%Figures: Electron-pair invariant mass after full selection, Fig. \ref{fig:InvMass}

 \begin{figure}[htb]
  \begin{center}
    \includegraphics[width=360pt]{Figures/invMass-04Apr11.png}
  \end{center}
  \caption[Electron-pair invariant mass after full selection]{Electron-pair invariant mass after full selection.}
  \label{fig:InvMass}
 \end{figure}

% Also add plots for OS and SS??  Had them at one point!  
% WELL, last OS plot I have does not show good data/MC agreement

%% Figures: Same-sign dielectron invariant mass for electrons passing all selection steps, Fig. \ref{fig:SameSignInvMass}

%%  \begin{figure}[htb]
%%   \begin{center}
%%     \includegraphics[width=360pt]{CMS-BW.pdf}
%%   \end{center}
%%   \caption[Same-sign dielectron invariant mass for electrons passing all selection steps]{Same-sign dielectron invariant mass for electrons passing all selection steps.}
%%   \label{fig:SameSignInvMass}
%%  \end{figure}

Ooh, I do also have inv mass by eta plots

The total number of data events surviving 
the selection cuts is 8453.  
% mention slight upward shift in data of peak position here?

ALSO TALK ABOUT HOW PLOTS SCALED??!!


recap of formula:

THE FORMULA

   * ntotal
   * nbg -- THE DIFFERENT METHODS OF BACKGROUND SUBTRACTION
   * acc (event sel)
   * eff (event sel) -- but talk about correcting for data/MC difference
   * Lumi (exp chapter or sth)

ntotal comes from inv mass plot just shown, nbg is what you need to calculate.
3 methods: 
   * MC estimate
   * template method
   * sideband subtraction

You know, I don't think my analysis needs a whole giant section on data vs MC comparison.  
I didn't do a whole lot.  
Well, just write up what I did do, I guess.  

%\section{Comparison of Monte Carlo and Data}
\section{Background Subtraction}
It is expected that some number of non-signal events %, 
%or background events, 
will look enough like signal events to pass the full selection.  
These background events need to be accounted for 
in the final calculation to arrive 
at the most accurate result.  
Therefore, the number of background events should 
be estimated and subtracted from the total number 
of events; 
This estimation can be done in a variety of ways.  
The methods used in this analysis are detailed 
in the following sections.  

Figures: Background subtracted yield versus eta with a MZ mass window - comparison to MC, Fig. \ref{fig:YieldVsZEtaComparedToMC}

 \begin{figure}[htb]
  \begin{center}
    \includegraphics[width=360pt]{CMS-BW.pdf}
  \end{center}
  \caption[Background subtracted yield versus eta with a MZ mass window - comparison to MC]{Background subtracted yield versus eta with a MZ mass window - comparison to MC.}
  \label{fig:YieldVsZEtaComparedToMC}
 \end{figure}

Figures: Background subtracted Z PT distribution with a MZ mass window - comparison to MC, Fig. \ref{fig:YieldVsZPtComparedToMC}

 \begin{figure}[htb]
  \begin{center}
    \includegraphics[width=360pt]{CMS-BW.pdf}
  \end{center}
  \caption[Background subtracted Z PT distribution with a MZ mass window - comparison to MC]{Background subtracted Z PT distribution with a MZ mass window - comparison to MC.}
  \label{fig:YieldVsZPtComparedToMC}
 \end{figure}

\subsection{Monte Carlo Estimation}
One estimate of the number of background events can be made 
directly from the Monte Carlo simulation, 
by counting the number of simulated events that pass 
the full selection.  
The numbers of passing events for the simulated background 
samples used are in Table FIXME, 
with a total of 18.6 events.  
The accuracy of this method may be questionable due to 
unaccounted differences bewteen the simulation and real data, however.  
In particular, 
the detector may not be correctly modeled in the simulation, 
causing slight differences in event numbers 
that are non-negligible at this level.  
Therefore, data-driven methods of background subtraction 
are also used, 
and the results from all the methods are checked against each other 
to confirm the accuracy of each method.  

TABLE with MC-estimated number of events from each background source

\subsection{Template Method}


\subsection{Sideband Subtraction}



\section{Estimation of Systematic Uncertainties}
MAYBE MAKE WHOLE OTHER CHAPTER???

Woohoo, the whole thing

   * lumi

   * stat (but that's not syst)

   * theory systs (HOW TO REFERENCE??  Just reference VBTF paper, I guess...?)

my stuff: 

   * electron energy scale
  
   * MC sample for efficiency

   * efficiency fitting

   * background subtraction/modeling (variation among the methods)

\section{What Else?}
Is this supposed to be random collection of methods used?  
e.g. tag and probe, template method, sideband fitting... 
although I guess the last two are actually background subtraction, 
which is somehow already determined to be part of this chapter.  